\startsection[title={Coleta de Dados}]
	A coleta das variáveis temperatura, potencial de oxirredução, concentração de hidrogênio, concentração de metano e concentração de dióxido de carbono é realizada {\emph on-line}. Dessa forma, a sincronização das medições é gerenciada automaticamente pelo controlador.

	A temperatura é controlada com {\emph set point} de \math{80 \celsius}, enquanto as demais variáveis são apenas monitoradas. Essa configuração decorre de limitações operacionais do sistema, que impedem a utilização de agitação mecânica e o controle do \chemical{pH}. Assim, uma das propostas deste trabalho é desenvolver um \getbuffer[sens] capaz de estimar variáveis de qualidade e de concentração em condições de operação sem agitação e sem controle de \chemical{pH}.
\stopsection

\startsection[title={Simulação}]
	Além da coleta de dados experimentais, também são gerados dados por simulação. Para isso, utiliza-se o modelo ADM1 acrescido de ruído gaussiano aditivo com variância compatível com a instrumentação experimental, introduzidas de forma controlada para aproximar os resultados do comportamento observado em sistemas reais e avaliar a capacidade da rede neural em identificar padrões na presença de ruído.

	O modelo adotado corresponde à variante proposta por \cite[authoryears][ref::rosen_0], incluindo as aproximações por equações diferenciais algébricas para o cálculo do \chemical{pH} e da concentração de \chemical{H_2} descritas neste trabalho.
\stopsection

\startsection[title={Arquitetura}]
	Os dados coletados pelos sensores e armazenados no servidor passam inicialmente por uma etapa de pré-processamento. Nessa etapa, são removidos valores inconsistentes e registros incompletos, como instantes em que alguma variável medida esteja ausente. Também é realizada uma inspeção visual das séries temporais para identificar possíveis anomalias.

	Em seguida, os dados são normalizados para adequá-los ao treinamento da rede neural e reduzir problemas numéricos, como instabilidades decorrentes de diferenças de escala entre as variáveis. O modelo LSTM é então treinado tanto com os dados experimentais quanto com os dados simulados. Após o treinamento, seu desempenho é avaliado utilizando conjuntos de dados não empregados durante o treinamento. A \in{Figura}[fig:archtecture] mostra essa sequência de trabalho.

	\startplacefigure[title={Fluxograma da arquitetura de funcionamento da planta para aquisição de dados, treinamento e monitoramento}, reference={fig:archtecture}]
		\externalfigure[image/architecture.png][height=400pt]
	\stopplacefigure

	A principal aplicação de um \getbuffer[sens] consiste no monitoramento do processo em tempo real, permitindo detectar desvios das condições normais de operação e identificar possíveis falhas.

	À medida que o sistema fermentativo evoluir para uma plataforma totalmente automatizada, com controle de variáveis como temperatura, \chemical{pH} e alimentação em regime semi-batelada, a integração do \getbuffer[sens] com um sistema de detecção de falhas permitirá a implementação de estratégias de controle mais robustas. Esse sistema poderá atuar sobre a velocidade de agitação, a correção do \chemical{pH} por meio da adição de soluções ácidas ou básicas, o aquecimento ou resfriamento da jaqueta do fermentador, a alimentação de substrato durante o processo e a composição gasosa da atmosfera do reator, incluindo a concentração de \chemical{CO_{2}}.
\stopsection

